%%%PAGE 228 rot=0 legibility=good summary=正则动量(洛伦兹力冲量)与多层电磁场题、滑轮连接体机械能题、两镜多次反射、半圆形玻璃光路、半圆导线安培力与张力
%%%BLOCK id=228-01 ch=11 sec=洛伦兹力分解·动量定理 kind=method alt=none cont=none y=0.03-0.17
\textbf{1.} 正则动量 \quad $qBx=m\Delta V_y$ \quad \red{2018、桥梁公式考} \quad \red{25T}
%FIG p228_a 0.09 0.075 0.28 0.165
\notefig[0.25]{figures/p228_a.png}

$qB_1l_1+qB_2l_2=m\Delta V_y=m(V_0\cos\alpha+V_0\sin\alpha)$
%%%END
%%%BLOCK id=228-02 ch=11 sec=复合场·多层电磁场 kind=problem alt=09 cont=none y=0.17-0.40
%FIG p228_b 0.03 0.175 0.235 0.265
\notefig[0.25]{figures/p228_b.png}

\begin{problem}
①求在第$n$层时出B的角度$\theta$（与水平方向）

②粒子恰可从第$n$层过，但比荷大一点不能穿过。解释
\begin{solution}
① $nqBd=m\Delta V_y$

$\Delta V_y=V\sin\theta_n$

$nEqd=\dfrac12 mV_n^{2}$ \qquad $V_n=\sqrt{\dfrac{2nEqd}{m}}$

$\Delta V_y=\sqrt{\dfrac{2nEqd}{m}}\,\sin\theta_n$

$\therefore\sin\theta_n=\sqrt{\dfrac{nqBd}{2Em}}$ % CHECK: 按推导根号内分子应为 nqB^2 d，原稿疑漏平方

② $\sin\theta\le1$ \qquad $\dfrac{q}{m}\uparrow$ \qquad $d\downarrow$

$d\downarrow$ 故出不去
\end{solution}
\end{problem}
%%%END
%%%BLOCK id=228-03 ch=06 sec=机械能守恒·关联速度连接体 kind=problem alt=03,04 cont=none y=0.43-0.60
\textbf{2.}
%FIG p228_c 0.02 0.425 0.18 0.575
\notefig[0.25]{figures/p228_c.png}

$V_P=V_2=V_Q\cos\theta$ \qquad $\because 0<\theta\le90^\circ$ \quad $\therefore V_P<V_Q$

P先失重 $T<mg$ 后超重 $T>mg$ % CHECK: P质量为4m，此处"mg"疑为4mg；"<"">"处笔迹重描

故 $T\ne mg$ \quad 可等于 不恒等于

Q到最低点 $E_{Q\text{机}}$最大 $E_{P\text{机}}$最小

\[ 4mg\left(\frac{d}{\cos37^\circ}-d\right)=mgd\tan37^\circ+\frac12 mV_Q^{2}\qquad V_Q=\frac{\sqrt{2gd}}{2} \]
%%%END
%%%BLOCK id=228-04 ch=15 sec=反射·两镜多次反射 kind=note alt=none cont=none y=0.62-0.70
\textbf{3.}
%FIG p228_d 0.07 0.625 0.155 0.695
\notefig[0.2]{figures/p228_d.png}

\red{两\unclear{镜}区 多次反射问题}
%%%END
%%%BLOCK id=228-05 ch=15 sec=折射·半圆形玻璃砖 kind=note alt=none cont=none y=0.62-0.73
%FIG p228_e 0.47 0.62 0.595 0.725
\notefig[0.25]{figures/p228_e.png}

出入平行 必过最高点

$\dfrac{\sin i}{\sin r}=n$
%%%END
%%%BLOCK id=228-06 ch=11 sec=安培力·有效长度与张力 kind=note alt=02 cont=none y=0.73-0.85
%FIG p228_f 0.015 0.735 0.155 0.845
\notefig[0.25]{figures/p228_f.png}

$F_{\text{全}}=2BIL$

\underline{$F_{\text{半}}=\dfrac{F}{2}=BIL$} \struck{张力} \quad \red{上下张力合$=F=2BIL$} % 原稿 F半 式下有红色下划线
%%%END
%%%PAGEEND 228
